\subsubsection{逆元}

若
\[
  ab\equiv1\pmod m,
\]
则称 $b$ 为 $a$ 在模 $m$ 意义下的逆元，记作 $a^{-1}$。

\subsubsection{同余运算与约分}

\[
  a\equiv b\pmod m
  \iff a=km+b
  \iff m\mid(a-b).
\]
同余式两边可以同时加、减、乘同一个数，但除法不能随意进行。
若 $d=\gcd(k,m)$，则
\[
  ka\equiv kb\pmod m
  \iff
  a\equiv b\pmod{m/d}.
\]
在 $0,1,\ldots,m-1$ 中，方程
\[
  Ax\equiv0\pmod m
\]
恰有 $\gcd(A,m)$ 个解。

\subsubsection{中国剩余定理}

考虑同余方程组
\[
  \begin{cases}
    x\equiv a_1\pmod{m_1},\\
    x\equiv a_2\pmod{m_2},\\
    \cdots\\
    x\equiv a_k\pmod{m_k}.
  \end{cases}
\]
当模数 $m_1,\ldots,m_k$ 两两互质时，令
\[
  M=\prod_{i=1}^{k}m_i,
  \qquad M_i=\frac{M}{m_i}.
\]
计算 $M_i$ 在模 $m_i$ 下的逆元 $M_i^{-1}$，再令 $c_i=M_iM_i^{-1}$（此处不先取模）。
唯一解（模 $M$）为
\[
  x\equiv\sum_{i=1}^{k}a_ic_i
  \equiv\sum_{i=1}^{k}a_iM_i(M_i^{-1}\bmod m_i)
  \pmod M.
\]
若 $V=\max_i m_i$，逐个求逆并合并的常见复杂度写作 $O(k\log V)$。

\subsubsection{扩展中国剩余定理}

模数不互质时，逐个合并方程。
先考虑
\[
  x\equiv a_1\pmod{m_1},
  \qquad
  x\equiv a_2\pmod{m_2}.
\]
由第一个方程写成 $x=a_1+m_1t$，代入第二个方程：
\[
  a_1+m_1t\equiv a_2\pmod{m_2},
\]
即
\[
  m_1t\equiv a_2-a_1\pmod{m_2}.
\]
它等价于一次不定方程
\[
  m_1t+m_2r=a_2-a_1.
\]
令 $g=\gcd(m_1,m_2)$，有解的充要条件为
\[
  g\mid(a_2-a_1).
\]
有解时约去 $g$：
\[
  \frac{m_1}{g}t
  \equiv
  \frac{a_2-a_1}{g}
  \pmod{m_2/g}.
\]
由于 $m_1/g$ 与 $m_2/g$ 互质，可以求逆元得到 $t$。
合并后的余数和模数为
\[
  a'=a_1+m_1t,
  \qquad
  m'=\operatorname{lcm}(m_1,m_2).
\]
依次两两合并即可。
若最终模数为 $N=\operatorname{lcm}(m_1,\ldots,m_k)$，逐个合并的常见复杂度写作 $O(k\log N)$。

\paragraph{应用：拆模数计算}

若要求答案模合数 $M$，先分解
\[
  M=p_1^{q_1}p_2^{q_2}\cdots p_s^{q_s}.
\]
分别计算
\[
  \mathrm{ans}\bmod p_1^{q_1},
  \quad\ldots,\quad
  \mathrm{ans}\bmod p_s^{q_s},
\]
最后用 CRT 合并成 $\mathrm{ans}\bmod M$。

\begin{icpcwarning}[实现细节]
  中间乘法可能溢出，按数据范围使用 \texttt{\_\_int128} 或快速乘。
  C++ 的负数取余可能为负，统一规范化到 $[0,m)$。
\end{icpcwarning}
